\section{Integers and Overflow}\label{sec:integer-overflows}

Discrete quantities are represented as binary signed or unsigned integers,
which can overflow when their values exceed the maximum or minimum representable values.
The range of an $n$-bit signed integer is $[-2^{n-1}, 2^{n-1}-1]$,
while an unsigned integer has a range of $[0, 2^n-1]$.

When an integer exceeds its maximum value, it wraps around to the minimum value (and vice versa).
Therefore, it is crucial to use sufficiently large integer types and/or implement overflow checking mechanisms.
This requires careful consideration of expected number ranges.

